\documentclass[11pt]{article}
\usepackage[a4paper,margin=1in]{geometry}
\usepackage{amsmath,amssymb,amsthm}
\usepackage[T1]{fontenc}
\usepackage{lmodern}
\usepackage[hidelinks]{hyperref}
\setlength{\emergencystretch}{3em}

\theoremstyle{plain}
\newtheorem{theorem}{Theorem}
\newtheorem{lemma}[theorem]{Lemma}
\newtheorem{corollary}[theorem]{Corollary}
\theoremstyle{definition}
\newtheorem{definition}[theorem]{Definition}
\theoremstyle{remark}
\newtheorem{remark}[theorem]{Remark}

\title{A Counterexample to Conjecture 00000002042:\\
Exact Covering Systems with Distinct Moduli Do Not Exist}
\author{%
  Feiyu\thanks{Independent researcher.}
  \and
  Claude (Opus 5.5)\thanks{Anthropic.}%
}
\date{2026-10-02}

\begin{document}
\maketitle

\begin{abstract}
Conjecture 00000002042 asserts that the number of exact covering systems with
distinct moduli grows like $c\,N^{1/2}$. By the theorem of Mirsky and Newman,
the only such system is the trivial one $\{0\bmod 1\}$, so the count is at most
$1$ for every $N$. We give the roots-of-unity proof and formalize it, together
with the refutation, in Lean~4 against Mathlib.
\end{abstract}

\section{The conjecture as stated}

The original text of conjecture 00000002042 reads:

\begin{quote}
Definition: An exact covering system covers every integer by exactly one
congruence class. Conjecture: The count of exact covering systems with distinct
moduli is asymptotically $c\cdot N^{1/2}$ (with $c$ of the explicit $2/\pi$
type); the count is controlled by the enumeration of Zn\'am-type trees.
\end{quote}

\section{Reading of the statement}

\begin{definition}
An \emph{exact covering system} is a finite set of residue classes $a_i\bmod
m_i$ ($m_i\ge1$, $0\le a_i<m_i$) such that every integer lies in exactly one of
them. It has \emph{distinct moduli} if the $m_i$ are pairwise distinct.
\end{definition}

Let $\mathrm{count}(N)$ be the number of exact covering systems with distinct
moduli, all at most $N$. The conjecture asserts that
$\mathrm{count}(N)/(c\sqrt N)\to1$ for some $c>0$. Other ways of measuring the
size $N$ (largest modulus, least common multiple) lead to the same refutation,
since by Theorem~\ref{thm:mn} there is at most one system in total.

\section{Result}

\begin{theorem}[Mirsky--Newman]\label{thm:mn}
The only exact covering system with distinct moduli is $\{0\bmod1\}$.
\end{theorem}

\begin{corollary}
$\mathrm{count}(N)\le1$ for all $N$, so $\mathrm{count}(N)/(c\sqrt N)\to0$ for
every $c>0$, and Conjecture 00000002042 is false.
\end{corollary}

\section{Proof}

\begin{proof}[Proof of Theorem~\ref{thm:mn}]
Let $M$ be the largest modulus, attained by the class $a^\ast\bmod M$, and
suppose $M\ge2$. Let $\zeta=e^{2\pi i/M}$, a primitive $M$-th root of unity, and
let $L=\prod_i m_i$, so every modulus divides $L$. Since $\zeta\ne1$ and
$\zeta^L=1$,
\[
  \sum_{n=0}^{L-1}\zeta^n=0 .
\]
Each $n\in[0,L)$ lies in exactly one class, and the elements of $[0,L)$ in the
class $a\bmod m$ are $a+mj$, $0\le j<L/m$. Hence
\[
  0=\sum_i\ \zeta^{a_i}\sum_{j=0}^{L/m_i-1}\bigl(\zeta^{m_i}\bigr)^j .
\]
If $m_i<M$, then $\zeta^{m_i}\ne1$ while $(\zeta^{m_i})^{L/m_i}=\zeta^L=1$, so the
inner geometric sum vanishes. Because the moduli are distinct, only the class
$a^\ast\bmod M$ remains, and it contributes $\zeta^{a^\ast}\cdot L/M\ne0$. This
is a contradiction, so $M=1$. Then every modulus is $1$, distinctness leaves a
single class, and it is $0\bmod1$.
\end{proof}

\section{Formalization}

The accompanying Lean~4 project (\texttt{lean/Submission/Basic.lean}) states
the conjecture and proves its negation against Mathlib.

\begin{itemize}
  \item \texttt{IsExactCovering S}, \texttt{DistinctModuli S} --- for a finite
    set \texttt{S} of pairs $(m,a)$ with $0<m$, $a<m$: every integer is
    congruent mod $m$ to $a$ for exactly one pair; the moduli are pairwise
    distinct.
  \item \texttt{count N} --- the number of such systems with all moduli at most
    $N$; \texttt{ConjectureHolds} --- $\exists c>0$,
    $\mathrm{count}(N)/(c\sqrt N)\to1$.
  \item \texttt{sum\_filter\_mod} --- the re-indexing of a residue class in
    $[0,L)$; \texttt{eq\_trivial} --- Theorem~\ref{thm:mn}, with $\zeta$ from
    Mathlib's \texttt{Complex.isPrimitiveRoot\_exp} and the geometric sums
    from \texttt{geom\_sum\_eq}.
  \item \texttt{count\_le\_one}, \texttt{conjecture\_00000002042\_false}.
\end{itemize}

The toolchain is \texttt{leanprover/lean4:v4.33.1}, Mathlib is pinned to
\texttt{v4.33.1}, and \texttt{lake build} succeeds. The project contains no
\texttt{sorry}, no \texttt{admit} and no \texttt{native\_decide}, and
introduces no axioms:
\begin{center}
\texttt{\#print axioms conjecture\_00000002042\_false}
\end{center}
reports exactly \texttt{propext}, \texttt{Classical.choice} and
\texttt{Quot.sound}.

\section{Remarks}

Theorem~\ref{thm:mn} is due to Mirsky and Newman, and independently to
Davenport and Rado, answering a question of Erd\H{o}s; see \cite{guy}. Exact
covering systems \emph{with repeated moduli} exist in abundance (for instance
$0\bmod2$, $1\bmod4$, $3\bmod4$), and the Zn\'am-type constructions mentioned
in the statement concern those. With distinct moduli the count is identically
$1$.

\begin{thebibliography}{9}
\bibitem{guy} R.~K.~Guy, \emph{Unsolved Problems in Number Theory}, 3rd ed.,
  Springer, 2004, Section F13.
\end{thebibliography}

\end{document}
